\appendix

\section{Hyper-Parameters}

Hyper-parameters related to the model architecture of GLM-5 are shown in \Cref{tab:model_arch_glm45_glm5}.

For training, we follow the setting of GLM-4.5, including the Muon optimizer, cosine decay, and batch size warmup. The learning rate goes through a warmup stage from 0 to 2e-4, and a decaying stage to 4e-5 until the end of the pre-training stage. In the mid-training stage, the learning rate decreases linearly from 4e-5 to 1e-5. Other hyper-parameters are the same as those of GLM-4.5. For DSA warmup stage, the learning rate goes down from 5e-3 to 2e-4. For DSA sparse adaption stage, we use a constant learning rate of 1e-5.

\begin{table}[!ht]
    \centering
    \caption{Model architecture of GLM-4.5 and GLM-5. When counting parameters, for all models we include the parameters of MTP layers but not word embeddings and the output layer.}
    \label{tab:model_arch_glm45_glm5}
    \begin{tabular}{l|cc}
    \toprule
       Model & \textbf{GLM-4.5} & \textbf{GLM-5} \\
    \midrule
       \# Total Parameters & 355B & 744B \\
       \# Activated Parameters & 32B & 40B \\
       \# Dense Layers & 3 & 3 \\
       \# MoE Layers & 89 & 75 \\
       \# MTP Layers & 1 & 1 \\
       Hidden Dim & 5120 & 6144 \\
       Dense Intermediate Dim & 12288 & 12288 \\
       MoE Intermediate Dim & 1536 & 2048 \\
    \midrule
       QK Head Dim & 128 & 192 \\
       V Head Dim & 128 & 256 \\
       Q LoRA Dim & -- & 2048 \\
       KV LoRA Dim & -- & 512 \\
       \# Attention Heads & 96 & 64 \\
       \# Key-Value Heads & 8 & -- \\
       \# Indexer Attn Heads & -- & 32 \\
       \# Indexer Head Dim & -- & 128 \\
    \midrule
       \# Experts (total) & 160 & 256 \\
       \# Routed Experts & 8 & 8 \\
       \# Shared Experts & 1 & 1 \\
    \midrule
       Vocabulary Size & 151552 & 154880 \\
    \bottomrule
    \end{tabular}
\end{table}

\section{Evaluation Details}
\subsection{Evaluation of Base Models}
We evaluate the base model of GLM-5 with English, Chinese, code, and math benchmarks in \Cref{tab:base_benchmark}.

\begin{table}[!ht]
  \centering
  \caption{Comparison among GLM-5-Base, GLM-4.5-Base and other representative open-source base models.}
  \resizebox{\textwidth}{!}{
  \begin{tabular}{lcccc|c}
    \toprule
     & \textbf{Benchmark (Metric)} & \textbf{DeepSeek-V3} & \textbf{Kimi-K2} & \textbf{GLM-4.5} & \textbf{GLM-5} \\
    & & \textbf{Base} & \textbf{Base} & \textbf{Base} & \textbf{Base} \\
    \midrule
    \multirow{3}{*}{} & Architecture & MoE & MoE & MoE & MoE \\
    & \# Activated Params & 37B & 32B & 32B & 40B \\
    & \# Total Params & 671B & 1043B & 355B & 744B \\
    \midrule
    \multirow{6}{*}{\textbf{English}} & SimpleQA (EM) & 26.6 & 35.3 & 30.0 & 36.0 \\
    & BBH (EM) & 88.4 & 88.7 & 86.2 & 87.4 \\
    & MMLU (EM) & 87.2 & 87.8 & 86.1 & 88.3 \\
    & HellaSwag (EM) & 88.9 & 94.6 & 87.1 & 88.1 \\
    & PIQA (EM) & 84.7 & - & 85.3 & 84.6 \\
    & TriviaQA (EM) & 82.9 & 85.1 & 80.0 & 80.9 \\
    \midrule
    \multirow{2}{*}{\textbf{Code}} & EvalPlus (Pass@1) & 65.6 & 80.3 & 78.1 & 87.0 \\
    & LiveCodeBench-Base (Pass@1) & 24.6 & 26.3 & 28.1 & 34.4 \\
    \midrule
    \multirow{2}{*}{\textbf{Math}} & GSM8K (EM) & 87.6 & 92.1 & 79.4 & 68.8 \\
    & MATH (EM) & 62.6 & 70.2 & 61.0 & 56.4 \\
    \midrule
    \multirow{4}{*}{\textbf{Chinese}} & CLUEWSC (EM) & 82.7 & - & 83.5 & 84.2 \\
    & C-Eval (EM) & 90.1 & 92.5 & 86.9 & 88.8 \\
    & C3 (EM) & 78.6 & - & 83.1 & 80.3 \\
    & Chinese-SimpleQA (EM) & 72.1 & 77.6 & 70.1 & 74.6 \\
    \bottomrule
  \end{tabular}
  }
  \label{tab:base_benchmark}
\end{table}

\subsection{Evaluation of ARC Benchmarks}
\label{sec:eval_details_arc}
Humanity’s Last Exam (HLE) \& other reasoning tasks: We evaluate with a maximum generation length of $131,072$ tokens ($temperature=1.0, top\_p=0.95, max\_new\_tokens=131072$). By default, we report the text-only subset; results marked with * are from the full set. We use GPT-5.2 (medium) as the judge model. For HLE-with-tools, we use a maximum context length of $202,752$ tokens.

SWE-bench \& SWE-bench Multilingual: We run the SWE-bench suite with OpenHands using a tailored instruction prompt. Settings: $temperature=0.7, top\_p=0.95, max\_new\_tokens=16384$, with a 200K context window.

BrowseComp: Without context management, we retain details from the most recent 5 turns. With context management, we use the same discard-all strategy as DeepSeek-V3.2 and Kimi K2.5.

Terminal-Bench 2.0 (Terminus 2): We evaluate with the Terminus framework using $timeout=2h, temperature=0.7, top\_p=1.0, max\_new\_tokens=8192$, with a 128K context window. Resource limits are capped at 16 CPUs and 32 GB RAM.

Terminal-Bench 2.0 (Claude Code): We evaluate in Claude Code 2.1.14 (think mode) with $temperature=1.0, top\_p=0.95, max\_new\_tokens=65536$. We remove wall-clock time limits, while preserving per-task CPU and memory constraints. We fix environment issues introduced by Claude Code and also report results on a verified Terminal-Bench 2.0 dataset that resolves ambiguous instructions (see: \url{https://huggingface.co/datasets/zai-org/terminal-bench-2-verified}). Scores are averaged over 5 runs.

CyberGym: We evaluate in Claude Code 2.1.18 (think mode, no web tools) with ($temperature=1.0, top\_p=1.0, max\_new\_tokens=32000$) and a 250-minute timeout per task. Results are single-run Pass@1 over 1,507 tasks.

MCP-Atlas: All models are evaluated in think mode on the 500-task public subset with a 10-minute timeout per task. We use Gemini 3 Pro as the judge model.

$\tau^2$-Bench: We add a small prompt adjustment in Retail and Telecom to avoid failures caused by premature user termination. For Airline, we apply the domain fixes proposed in the Claude Opus 4.5 system card.

Vending-Bench 2: Runs are conducted independently by Andon Labs\footnote{\url{https://andonlabs.com/evals/vending-bench-2}}.


\subsection{Optimized User Simulator for $\tau^2$-Bench}
\label{sec:tau_user_prompt}
We add a small prompt adjustment in Telecom and Retail to avoid failures caused by premature user termination. The optimized prompts are shown in Figure~\ref{fig:tau2-prompt-telecom} and Figure~\ref{fig:tau2-prompt-retail}. These optimized prompts are integrated into the system prompt as follows:

\begin{tcolorbox}[left=0mm,right=0mm,top=0mm,bottom=0mm,boxsep=1mm,arc=0mm,boxrule=0pt, frame empty, breakable]
\begin{lstlisting}
SYSTEM_PROMPT = """"
{global_user_sim_guidelines}

<scenario>
{instructions}
</scenario>

{optimized_user_prompt}
"""".strip()
\end{lstlisting}
\end{tcolorbox}

\begin{tcolorbox}[left=0mm,right=0mm,top=0mm,bottom=0mm,boxsep=1mm,arc=0mm,boxrule=0pt, frame empty, breakable]
\footnotesize
\begin{lstlisting}
# Note:
- Do not generate the '###TRANSFER###' before agent clearly tells "YOU ARE BEING TRANSFERRED TO A HUMAN AGENT. PLEASE HOLD ON.".
    Example:
    Case1:
        - agent: "Would you like me to transfer you to a human agent who can assist you with these options and help get your service restored?"
        - user: "Yes, please transfer me to a human agent.".
    Case2:
        - user: "YOU ARE BEING TRANSFERRED TO A HUMAN AGENT. PLEASE HOLD ON."
        - user: "###TRANSFER###"
\end{lstlisting}
\end{tcolorbox}
\noindent\begin{minipage}{\textwidth}
\captionof{figure}{The optimized user prompt for $\tau^2$-Bench Telecom.}
\label{fig:tau2-prompt-telecom}
\end{minipage}

\begin{tcolorbox}[left=0mm,right=0mm,top=0mm,bottom=0mm,boxsep=1mm,arc=0mm,boxrule=0pt, frame empty, breakable]
\footnotesize
\begin{lstlisting}
# Rules:
- Just generate one line at a time to simulate the user's message.
- Do not give away all the instruction at once. Only provide the information that is necessary for the current step.
- Do not hallucinate information that is not provided in the instruction. Follow these guidelines:
    1. If the agent asks for information NOT in the instruction:
        - Say you don't remember or don't have it
        - Offer alternative information that IS mentioned in the instruction
    2. Examples:
        - If asked for order ID (not in instruction): "Sorry, I don't remember the order ID, can you search for it? My name/email/phone number/zipcode is ..."
        - If asked for email (not in instruction): "I don't have my email handy, but I can give you my name and zip code which are..."
- Do not repeat the exact instruction in the conversation. Instead, use your own words to convey the same information.
- Try to make the conversation as natural as possible, and stick to the personalities in the instruction.
# Constraint Handling:
- Provide requests strictly based on what is explicitly stated in the instruction.
- Do not assume, extend, substitute, or generalize in any form.
- Do not modify or relax constraints on:
- Time / Date
- Budget
- Specific terms (e.g., "same" must not be replaced with "similar")
- Core Rule: Any attribute NOT mentioned in the instruction can be either changed or kept the same
- Examples:
    - If instruction says "exchange red item to blue": Only color must change, other attributes (size, material, etc.) are flexible
    - If instruction says "exchange red item to blue, keep the same size": Both color must change AND size must stay the same
- Exception: Only follow additional constraints when explicitly stated in the instruction
# Domain-Specific Rules:
## For Retail scenarios:
- Focus on product attributes and exchange/return processes as specified in instructions.
- During confirmations: Always respond based strictly on the original instruction, never deviate to match agent's provided options. Restate your requirement from the instruction rather than selecting from agent's choices.
    - Example: If the agent provides specific options (A/B/C) but the instruction states a general requirement (e.g., "same as pending order"), always restate or confirm based on what the instruction says, not by directly selecting from the agent's provided options.
# When NOT to finish the conversation:
- Do not end until you have clearly and completely expressed all your requirements and constraints.
- Do not end until the agent has completed all tasks mentioned in the instruction and verified no operations were missed.
- Do not end if the agent's execution results do not match your expectations or are incorrect/incomplete.
# When you CAN finish the conversation:
- Only when all above conditions are satisfied AND all tasks are completed correctly.
- OR when you have clearly expressed complete requirements but the system explicitly states it cannot complete them due to technical limitations - in this case, accept transfer to human.
# How to finish the conversation:
- If the agent has completed all tasks, generate the '###STOP###' token to end the conversation.
# Note:
- You should carefully check if the agent has completed all tasks mentioned in the instruction before generating '###STOP###'.
\end{lstlisting}
\end{tcolorbox}
\noindent\begin{minipage}{\textwidth}
\captionof{figure}{The optimized user prompt for $\tau^2$-Bench Retail.}
\label{fig:tau2-prompt-retail}
\end{minipage}

\subsection{Evaluation of Real-world Agentic Engineering Experience}

\subsubsection{Frontend Evaluation}

\label{sec:frontend_eval_data_cons}

\paragraph{Data.}
Our dataset encompasses seven distinct frontend scenarios designed to evaluate a model's engineering proficiency across diverse functional domains: Business Management Systems, Web Games, SVG/Canvas Rendering,  Creative Tools \& Editors, Showcase Pages, Forms \& Tables and Data Visualization.
% \begin{itemize}
% \item Business Management Systems: Applications for enterprise or personal data/workflow management.
% \item Web Games: Interactive, real-time rendering-focused entertainment applications.
% \item SVG/Canvas Rendering: Graphics-intensive applications based on low-level drawing APIs.
% \item Creative Tools \& Editors: Utility applications supporting content creation and complex state management.
% \item Showcase Pages: Layout-oriented pages emphasizing visual expression and responsiveness.
% \item Forms \& Tables: Applications centered on structured data entry, validation, and processing.
% \item Data Visualization: Graphical representation and analytical tools for complex datasets.
% \end{itemize}


\begin{table}[h]
\centering
\small
\caption{Distribution of frontend application scenarios.}
\begin{tabular}{lp{7cm}rr}
\toprule
\textbf{Category} & \textbf{Description} & \textbf{\# Tasks} & \textbf{\# Checkitems} \\
\midrule
Business Systems & Enterprise/Personal data and process management. & 42 & 167 \\
Web Games & Interaction and entertainment-focused games. & 40 & 163\\
SVG/Canvas & Graphics rendering and interactive visualizations. & 32 & 166\\
Creative Tools & Content creation and online editing tools. & 28 & 160\\
Showcase Pages & Visual expression and information presentation. & 27 & 115\\
Forms \& Tables & Structured data entry and processing. & 26 & 93\\
Data Visualization & Graphical data expression and analysis. & 25 & 85\\
\bottomrule
\end{tabular}

\end{table}

\subparagraph{Data Distribution by Coding Languages}
The benchmark provides full coverage of three mainstream paradigms: Vanilla Web Stack (HTML/CSS/JS), React Component-based Framework, and the Vue 3 + Vite Progressive Solution.
\begin{table}[h]
\centering
\caption{Statistics of technology stacks and evaluation units.}
\begin{tabular}{llcc}
\toprule
\textbf{Category} & \textbf{Description} & \textbf{\# Tasks} & \textbf{\# Checkitems} \\
\midrule
HTML & Vanilla HTML/CSS/JS development. & 113 & 490 \\
React & Component-based framework development. & 58 & 249 \\
Vue & Vue 3 + Vite progressive solution. & 49 & 210 \\
\bottomrule
\end{tabular}
\end{table}

\subparagraph{Data sample}
Each test case is composed of three components: the \textbf{Task}, the \textbf{Checklist}, and a \textbf{Dedicated Environment}. Below is a representative example of a test case:
\begin{tcolorbox}[left=0mm,right=0mm,top=0mm,bottom=0mm,boxsep=1mm,arc=0mm,boxrule=0pt, frame empty, breakable]
\footnotesize
\begin{lstlisting}
    Task: Develop an online drawing tool that includes a brush, an eraser, a white canvas, and a save button.
        The brush color and thickness should be selectable via buttons on the left. Users can draw on the canvas by clicking and dragging the mouse.
        The eraser size should be selectable via buttons on the left. Users can erase content by clicking and dragging the mouse over the canvas.
        Once the drawing is complete, clicking the "Save" button should allow the user to save the image locally.
        Please implement this using the React framework in the current directory.

    Checklist:
        The user can select the brush color and thickness using the left-hand buttons, and drawing is functional via mouse click-and-drag on the canvas.
        The user can select the eraser size using the left-hand buttons, and erasing is functional via mouse click-and-drag on the canvas.
        Upon clicking the "Save" button, the generated image is successfully saved to the local machine.
\end{lstlisting}
\end{tcolorbox}
\noindent\begin{minipage}{\textwidth}
\end{minipage}


\subparagraph{\textbf{Data Construction and Validation}}

We implement a rigorous four-stage pipeline to ensure data quality:
\begin{itemize}[leftmargin=1.5em,itemsep=2pt]
\item Stage 1: Task Synthesis. Tasks are designed by senior frontend experts to ensure they reflect real-world engineering challenges while maintaining a balanced distribution across diverse scenarios and technologies.
\item Stage 2: Checklist Generation and Refinement. We initially employ Claude Sonnet 4.5 to synthesize candidate checklists based on task specifications $T$. These are then meticulously audited and integrated by experts. Through multiple rounds of refinement, we ensure that each check-item is semantically unambiguous, objective, and provides exhaustive coverage of user requirements.
\item Stage 3: Execution-based Correction. We conduct cross-validation between the Agent-as-a-Judge framework and human experts. Any discrepancies in judgment trigger a re-evaluation and correction of the underlying data to eliminate potential noise.
\item Stage 4: Dynamic Benchmark Iteration. To maintain a high level of discriminative power, we iteratively update the test suite by removing trivial tasks that no longer challenge state-of-the-art coding agents. This expert-led curation process culminated in a final set of 220 high-quality frontend coding tasks and their corresponding checklists.
\end{itemize}